% Generated draft: requires mathematical and source review.
\documentclass[11pt]{article}
\usepackage[margin=25mm]{geometry}
\usepackage{fontspec}
\setmainfont{Latin Modern Roman}
\usepackage{amsmath,amssymb,mathtools,unicode-math}
\setmathfont{Latin Modern Math}
\usepackage{enumitem,needspace,xurl,hyperref}
\setlength{\parindent}{0pt}
\setlength{\parskip}{5pt}
\setlength{\emergencystretch}{3em}
\newcommand{\N}{\mathbb{N}}
\newcommand{\Z}{\mathbb{Z}}
\newcommand{\Q}{\mathbb{Q}}
\newcommand{\R}{\mathbb{R}}
\newcommand{\C}{\mathbb{C}}
\newcommand{\F}{\mathbb{F}}
\newcommand{\A}{\mathbb{A}}
\newcommand{\PP}{\mathbb{P}}
\newcommand{\E}{\mathbb{E}}
\newcommand{\eps}{\varepsilon}
\newcommand{\dd}{\,\mathrm{d}}
\begin{document}
\section*{UnsolvedMath ID: 30000838}
\textbf{Source problem code:} catalogue entry
\subsection{Definitions and mathematical statement}
A standard macroscopic setting for the question is a scalar conservation law
\[
\partial_t u+\nabla_x\!\cdot F(x,u)=0,\qquad u(0,x)=u_0(x),
\]
where the flux may jump across spatial interfaces. For example, on $\R^d$ a weak solution satisfies
\[
\int_0^\infty\!\int_{\R^d}
\bigl(u\,\partial_t\varphi+F(x,u)\cdot\nabla\varphi\bigr)\,\dd x\,\dd t
+\int_{\R^d}u_0(x)\varphi(0,x)\,\dd x=0
\]
for compactly supported smooth test functions. This weak equation alone does not specify admissibility at a flux discontinuity.

For the interacting particle system intended by the source, let $\pi_t^N$ be its empirical density measure after the prescribed space-time rescaling. The hydrodynamic selection problem is to identify an entropy or interface-admissibility condition $\mathcal E$ for which the macroscopic limit is the uniquely selected solution $u$ and
\[
\langle\pi_t^N,\varphi\rangle\longrightarrow
\int u(t,x)\varphi(x)\,\dd x
\]
in the intended probability and topology sense, for the admissible initial data and observation times. The task is to derive the physically relevant admissibility from the microscopic dynamics, including what it imposes at discontinuities, rather than merely postulating a weak solution.

The extract specifies neither the particle generator, rates or interaction class, nor the scaling, spatial domain, initial/boundary data, flux assumptions, convergence mode or entropy/interface condition. Even the variable in which the flux is discontinuous is not explicit; the displayed spatial-jump law is a conventional illustrative setting. Thus the record asks for a general selection principle, not a uniquely determined theorem for a supplied model. No adapted-entropy or vanishing-viscosity prescription is silently imposed for every microscopic model.
\subsection{Short English statement}
Determine which admissibility condition is forced by particle-system hydrodynamic limits for scalar laws with discontinuous flux. The microscopic process and interface rules are missing; known adapted-entropy results apply only to the supplied partial classes.
\subsection{Sources}
\begin{itemize}
\item Supplied report: \emph{Large Scale Stochastic Dynamics} (2008), Oberwolfach Reports, DOI 10.4171/owr/2007/42, pp. 2443--2530; proposed year 2007. \url{https://doi.org/10.4171/owr/2007/42}.
\item Supplied Chen--Even--Klingenberg reference: \url{https://www.math.ntnu.no/conservation/2007/026.html}.
\item Supplied Audusse--Perthame reference: \url{https://doi.org/10.1017/S0308210500003863}.
\item Local UnsolvedMath registry, permanent ID 30000838, catalogue entry; created 25 August 2026, updated 16 September 2026. Its statement is labelled exact and retained byte-for-byte. The supplied 21 August 2026 literature triage labels the record partially solved: existence, uniqueness and hydrodynamic selection of adapted entropy solutions are reported in important discontinuous-flux classes, while the question seeks a more general physically relevant selection principle. This dated assessment is not independently updated here.
\end{itemize}
\end{document}
